%!TEX root = ../../Diplomarbeit.tex
\section{Evaluation \& Discussion}

\subsection{Expected Engagement Improvements}
\begin{itemize}
    \item \textbf{Increased daily active users:} Gamification features such as streaks create a psychological incentive to log in every day, as users want to protect their accumulated progress.
    \item \textbf{Longer session duration:} Interactive elements and progression systems encourage users to explore additional features and spend more time within the app.
    \item \textbf{Better plant care consistency:} Reminder quests and rewards aim to regularize care activities, leading to healthier plants and reinforcing the connection between digital actions and real-world outcomes.
\end{itemize}

\subsection{Risks \& Limitations}
\begin{itemize}
    \item \textbf{Users gaming the system:} Users might perform unnecessary actions solely to earn points, such as over-watering a plant that does not need it.
    \item \textbf{Motivation drop after novelty:} The initial engagement spike from new features may fade over time as the ``novelty effect'' (the temporary boost in interest caused by something being new) wears off.
    \item \textbf{Balance issues:} Ensuring that rewards feel both earned and attainable requires careful calibration and ongoing adjustment.
\end{itemize}

\subsection{Comparison to Non-Gamified Apps}
\begin{itemize}
    \item \textbf{Traditional gardening apps:} These typically focus solely on providing information and scheduling reminders, without any engagement layer.
    \item \textbf{Why they lose users:} The absence of immediate feedback and an emotional hook (a design element that creates personal attachment to the app) leads to abandonment once initial enthusiasm fades or routine tasks begin to feel like chores.
\end{itemize}
